\chapter{因果中介}

第三部分介绍了估计时变处理对结局之因果效应的各种方法。我们的目标是量化结局分布在不同的、需长期维持的处理策略之下的变化。我们的目标并不是确定处理\emph{如何}对结局施加效应。现在我们把注意力转向因果中介：研究处理借以影响结局的因果通路。

因果中介的研究可以看作时变处理下因果推断的一个特例。中介分析中没有随时间取不同值的单一处理，而是在不同时点上拥有两个不同的变量——所关心的处理与中介。

本章描述因果中介的一个理论框架，它与本书其余部分一样，建立在可以映射为目标试验的假想干预之上。其他中介研究路径基于纯直接效应与总间接效应；干预主义的中介框架则不同，它为因果估计的经验验证打开了大门——这是任何科学事业都值得具备的条件。

\section{中介分析受到抨击}

考虑一项随机试验：从一个吸烟者总体中随机抽取样本，分配到戒烟（$A=0$）或继续吸烟（$A=1$）。假定依从性完美：所有被分配到 $A=0$ 的人都戒了烟，所有被分配到 $A=1$ 的人都继续吸烟。研究者发现戒烟对 1 年时心肌梗死风险 $Y$ 有有益效应，即 $\E[Y\given A=1]>\E[Y\given A=0]$。

戒烟 $A$ 对心脏病风险 $Y$ 的有益效应确立之后，研究者想知道这一益处是否源于 $A$ 降低高血压 $M$ 的效应。除基线处理 $A$ 与 1 年时结局 $Y$ 的数据外，他们还有 6 个月时测得的高血压 $M$ 的数据。（为简单起见，假定前 6 个月内没有人发生结局 $Y$。）在“对 $M$ 的干预定义得足够好”这一假定下，图 23.1 的因果图刻画了这些变量。

研究者感兴趣的是把 $A$ 对 $Y$ 的总效应分解为经由 $M$ 中介的间接因果通路与不经由 $M$ 中介的直接通路。此时我们说，研究者感兴趣的是\emph{因果中介分析}。

上一章介绍了几种把直接效应概念形式化的途径（技术注记 22.1 与 22.2），但没有介绍间接效应的概念。要把总处理效应分解为直接效应与间接效应，可以定义 $A$ 对 $Y$ 的纯直接效应与总间接效应。\mnote{Robins 与 Greenland (1992) 引入了纯直接效应与总间接效应。}

$A$ 对 $Y$ 的、不经过 $M$ 的\emph{纯直接效应}是：若对每个个体都把 $M$ 的取值设为其在 $A$ 被设为 0 时会取的取值，即把 $M$ 设为 $M^{a=0}$（对某些个体为 1，对另一些为 0），在此情形下 $A$ 对 $Y$ 的平均因果效应。于是纯直接效应就是对比
\[
\E\big[Y^{a=1,M^{a=0}}\big]-\E\big[Y^{a=0,M^{a=0}}\big].
\]

在我们的例子中，纯直接效应是：若我们能给每个人指派她在戒烟时会有的高血压取值 $M$（1 或 0），在此情形下戒烟 $A$ 的效应。对实际戒了烟的人，这个 $M$ 值已知；对继续吸烟的人，它仍然未知。因此我们说，$\E\big[Y^{a=1,M^{a=0}}\big]$、从而纯直接效应，是一个\emph{跨世界}量，因为它涉及一个由两个处理取值 $a=1$ 与 $a=0$ 索引的反事实结局，而这两个取值不可能在同一世界的同一个体身上同时出现。

\emph{总间接效应}同样是跨世界量，定义为
\[
\E\big[Y^{a=1,M^{a=1}}\big]-\E\big[Y^{a=1,M^{a=0}}\big].
\]

纯直接效应与总间接效应之和 $\E\big[Y^{a=1,M^{a=1}}\big]-\E\big[Y^{a=0,M^{a=0}}\big]$，由一致性，就是总效应 $\E[Y^{a=1}]-\E[Y^{a=0}]$。

在图 23.1 中，中介 $M$ 与结局 $Y$ 之间不存在未测量的共同原因。若存在共同原因，任何直接效应都无法识别，包括受控直接效应与中介效应，即纯直接效应与总间接效应。要识别中介效应，我们需要在图 23.1 的因果图下识别跨世界量 $\E\big[Y^{a=1,M^{a=0}}\big]$。所谓的\emph{中介公式}正是做这件事：
\[
\sum_m \E[Y\given A=1, M=m]\,\Pr[M=m\given A=0].
\]
证明见技术注记 23.1。

\begin{figure}[htbp]
\centering
\begin{tikzpicture}[causal]
  \node[vtx] (A) at (0,0) {$A$};
  \node[vtx] (M) at (2.2,0) {$M$};
  \node[vtx] (Y) at (4.4,0) {$Y$};
  \draw[ed] (A) -- (M);
  \draw[ed] (M) -- (Y);
  \draw[ed] (A) to[bend left=45] (Y);
\end{tikzpicture}
\caption{}
\label{fig:23.1}
\end{figure}

中介公式看似令人惊讶，因为它识别了一个跨世界反事实量，而这个量的取值无法通过任何对 $A$ 与 $M$ 的实验干预得到经验确证，哪怕在原则上也不行。然而，要做到这一点，需要假定图 23.1 的因果图表示的是 NPSEM-IE 而非 FFRCISTG；后者是本书通篇使用的反事实模型（见技术注记 6.2）。\mnote{在某些极不寻常的情形下，我们可以通过交叉试验识别这个跨世界反事实量的取值（见精要 2.1 与 3.2）。}

\begin{techpoint}[中介公式的证明]
\setlength{\parindent}{2em}
在图 23.1 的因果图所刻画的情境中（处理为 $A$、中介为 $M$、结局为 $Y$），反事实均值 $\E\big[Y^{a=1,M^{a=0}}\big]$ 等于
\begin{align*}
&= \sum_m \E\big[Y^{a=1,m}\given M^{a=0}=m\big]\,\Pr[M^{a=0}=m]
   &&\text{（由概率法则）}\\
&= \sum_m \E\big[Y^{a=1,m}\big]\,\Pr[M^{a=0}=m]
   &&\text{（由跨世界条件独立）}\\
&= \sum_m \E[Y\given A=1, M=m]\,\Pr[M=m\given A=0]
   &&\text{（由可交换性与一致性）}
\end{align*}
其中第二个等号用到的跨世界条件独立为 $Y^{a=1,m}\indep M^{a=0}\given M^{a=0}=m$。

上面用到的跨世界条件独立，在图 23.1 表示 NPSEM-IE 时被假定，在它表示 FFRCISTG 模型时则不然。
\end{techpoint}

中介公式的证明要求一个不可验证的假定：某些定义在不同世界中的反事实变量相互独立。具体而言，证明要求：在我们联合干预、把处理取值设为 1、把中介取值设为 $m$ 的世界里，结局 $Y^{a=1,m}$ 的取值，独立于我们把处理取值干预为 0 的世界里中介 $M^{a=0}$ 的取值。因此，变量 $Y^{a=1,m}$ 与 $M^{a=0}$ 不可能在同一世界的同一个体身上同时被观测到；于是在任何单独或联合随机分配 $A$ 与 $M$ 的试验中，它们的独立性都无法得到经验验证（见技术注记 7.1）。

这些\emph{跨世界独立性}在 NPSEM-IE 下被假定，在 FFRCISTG 模型下则不然。这是我们宁愿选择 FFRCISTG 而非 NPSEM-IE 的一个理由：既然任何随机试验都无法验证跨世界独立性，我们就希望使用其结果原则上可验证的因果推断方法。\mnote{在 FFRCISTG 模型下，纯直接效应与总间接效应不是点识别的，但可以得到锐界（Robins 与 Richardson 2010）。}

基于这一讨论，一些政策制定者对研究者能够识别纯直接效应并不买账。作为一个跨世界参数，他们认为纯直接效应没有政策或公共卫生意义，因为它不对应任何干预。要说服这些怀疑者，NPSEM-IE 的拥护者还得加把劲。

\section{为中介分析辩护}

戒烟试验的研究者作为 NPSEM-IE 的拥护者，准备了如下论证，以说服怀疑者相信纯直接效应的实际用途。

\begin{quote}
“假设从一年后开始，无尼古丁香烟将会上市，而你的政策目标是尽快了解无尼古丁香烟的益处。为此，我们用上述戒烟试验已收集的数据，估计在要求所有吸烟者改吸无尼古丁香烟（待其上市后）的干预下，两年期心脏病风险 $Y$。假设已有强有力的实验证据表明：(i) 尼古丁对结局 $Y$（心脏病）的全部因果效应都经由其对中介 $M$（高血压）的效应；(ii) 香烟中的非尼古丁成分对中介 $M$（高血压）没有直接因果效应。在假定 (i) 与 (ii) 之下，吸无尼古丁香烟者的高血压状态 $M$ 将等于她不暴露于香烟时的高血压状态，因此 $\E\big[Y^{a=1,M^{a=0}}\big]$ 就是所有吸烟者若在一年后改吸无尼古丁香烟时的心脏病风险。这正是你所关心的公共卫生效应，而中介公式应用于我们研究的数据时计算的正是这个量。”
\end{quote}

这很有意思。为了论证参数 $\E\big[Y^{a=1,M^{a=0}}\big]$ 的实质重要性，NPSEM-IE 的拥护者讲了一个关于干预香烟尼古丁含量的故事——而这个干预根本没有提到中介 $M$。此外，他们还作出了关于“无直接效应”的假定 (i) 与 (ii)，而这两个假定涉及的变量根本没有出现在图 23.1 的因果图上。\mnote{Pearl (2001) 提供了一个类似的论证，来为纯直接效应的使用辩护。}

研究者的故事声称，变量 $A$ 可以分解为两个可分离的成分 $N$ 与 $O$。可分离成分 $N$ 直接影响 $M$ 但不影响 $Y$。可分离成分 $O$ 直接影响 $Y$ 但不影响 $M$。此外，按照这个故事，每个可分离成分原则上都可以被单独干预。例如，$Y^{n=0,o=1}$ 表示在仅去除香烟中尼古丁成分的干预下的反事实结局。

这一因果故事最直接的表达，是图 23.2 的因果 DAG 所表示的 FFRCISTG 模型，其中 $N$ 是表示尼古丁暴露的二值变量，$O$ 是表示暴露于香烟其他非尼古丁成分的二值变量。从 $A$ 到 $N$ 与 $O$ 的粗箭头表示确定性关系。这是因为，在试验的实际数据中，要么继续吸普通香烟，于是 $A=N=O=1$；要么戒掉香烟，于是 $A=N=O=0$。$N$ 到 $Y$ 之间没有箭头，这编码了假定 (i)：$N$ 对 $Y$ 没有直接效应。$O$ 到 $M$ 之间没有箭头，这编码了假定 (ii)：$O$ 对 $M$ 没有效应。图 23.3 给出相应的 SWIG。若对每个个体 $O$ 都不是 $M$ 的原因，我们还可以把随机变量 $M^{n,o}$ 写作 $M^{n}$。\mnote{这些是“无受控直接效应”假定。假定 (i) 是说：在控制（设定）$M$ 的取值时，$N$ 对 $Y$ 没有直接效应。}

\begin{figure}[htbp]
\centering
\begin{tikzpicture}[causal]
  \node[vtx] (A) at (0,0) {$A$};
  \node[vtx] (N) at (2.0,1.1) {$N$};
  \node[vtx] (O) at (2.0,-1.1) {$O$};
  \node[vtx] (M) at (4.1,1.1) {$M$};
  \node[vtx] (Y) at (5.7,0) {$Y$};
  \draw[edb] (A) -- (N);
  \draw[edb] (A) -- (O);
  \draw[ed] (N) -- (M);
  \draw[ed] (M) -- (Y);
  \draw[ed] (O) -- (Y);
\end{tikzpicture}
\caption{}
\label{fig:23.2}
\end{figure}

\begin{figure}[htbp]
\centering
\begin{tikzpicture}[causal,swig/.style={draw,rounded corners=2pt,inner sep=2.5pt,minimum height=6mm,font=\small}]
  \node[vtx] (A) at (0,0) {$A$};
  \node[swig] (N) at (2.2,1.1) {$N\,|\,n$};
  \node[swig] (O) at (2.2,-1.1) {$O\,|\,o$};
  \node[swig] (M) at (4.4,1.1) {$M^{n,o}$};
  \node[swig] (Y) at (6.0,0) {$Y_{n,o}$};
  \draw[ed] (A) -- (N);
  \draw[ed] (A) -- (O);
  \draw[ed] (N) -- (M);
  \draw[ed] (M) -- (Y);
  \draw[ed] (O) -- (Y);
\end{tikzpicture}
\caption{}
\label{fig:23.3}
\end{figure}

在这一问题刻画下，识别 $\E\big[Y^{n=0,o=1}\big]$ 的 g 公式恰好就是中介公式（证明见技术注记 23.2）。下面回顾我们得到的教益。

\begin{techpoint}[当中介公式就是 g 公式时]
\setlength{\parindent}{2em}
按照图 23.3 的因果图，可交换性对可分离成分 $N$ 与 $O$ 成立，因此均值 $\E\big[Y^{n=0,o=1}\big]$ 可以由 g 公式识别。

然而，在正文描述的戒烟试验中，没有个体具有数据 $(N=0,O=1)$。因此正值性不成立，看来 g 公式无法识别 $\E\big[Y^{n=0,o=1}\big]$。但是，给定可交换性与一致性，正值性只是 g 公式识别的充分而非必要条件；识别只要求 g 公式是观测数据分布的函数。下面我们证明，图 23.2 中的确定性粗箭头与无直接效应假定 (i)、(ii) 合起来蕴含：g 公式确实是 $(A,M,Y)$ 的观测数据分布的函数——那个函数正是中介公式。

在图 23.2 所表示的 FFRCISTG 下，若有 $N$ 与 $O$ 的数据，均值 $\E\big[Y^{n=0,o=1}\big]$ 会由 g 公式
\[
\sum_m \E[Y\given O=1,M=m]\,\Pr(M=m\given N=0)
\]
识别，因为在 g 公式中 $N=0$ 不必列入第一个条件事件、$O=1$ 不必列入第二个条件事件，原因是在图 23.2 中 $N$ 不是 $Y$ 的父节点、$O$ 不是 $M$ 的父节点。此外，即使正值性不成立，该 g 公式也只是观测数据分布的函数，因为 $O=1$ 当且仅当 $A=1$，$N=0$ 当且仅当 $A=0$。于是我们可以把 $\E\big[Y^{n=0,o=1}\big]$ 的 g 公式改写为
\[
\sum_m \E[Y\given A=1,M=m]\,\Pr(M=m\given A=0),
\]
这正是中介公式。这一基于 g 公式的推导略显启发式，因为 $N$、$O$ 与 $A$ 之间的确定性会产生零测集。Robins 等 (2022) 用技术注记 21.12 所述的 SWIG 马尔可夫性质给出了另一种严格证明。

g 公式不仅在图 23.2 的扩展因果图下等于中介公式，而且在图 23.7 的扩展因果图下也等于前门公式，见技术注记 23.3。
\end{techpoint}

对一位最初接受与图 22.1 的 DAG 相联系的 NPSEM-IE 的研究者而言，均值 $\E\big[Y^{a=1,M^{a=0}}\big]$ 早已由中介公式识别，因此关于 $N$ 与 $O$ 之可分离效应的故事对识别毫无贡献。它只是用来说明：参数 $\E\big[Y^{a=1,M^{a=0}}\big]$——从而纯直接效应与总间接效应——编码了一个具有公共卫生意义的参数 $\E\big[Y^{n=0,o=1}\big]$。对持 FFRCISTG 观点的研究者而言，这个故事不仅把均值 $\E\big[Y^{a=1,M^{a=0}}\big]$ 解释为均值 $\E\big[Y^{n=0,o=1}\big]$，而且还使两个均值都能通过 g 公式、或等价地通过中介公式识别。事实上，当我们把“没有箭头”解释为“对所有个体都没有直接因果效应”时，$Y^{a=1,M^{a=0}}$ 与 $Y^{n=0,o=1}$ 对每个个体都相等。

这一刻画还使我们能够定义并识别成分 $N$ 对结局 $Y$ 的\emph{可分离直接效应} $\E\big[Y^{n=1,o=1}\big]-\E\big[Y^{n=0,o=1}\big]$，它是可分离成分 $N$ 的受控直接效应。这一效应也等于总间接效应 $\E[Y^{a=1}]-\E\big[Y^{a=1,M^{a=0}}\big]$。类似地，$\E\big[Y^{n=1,o=1}\big]-\E\big[Y^{n=1,o=0}\big]$ 是纯直接效应 $\E\big[Y^{a=1,M^{a=0}}\big]-\E[Y^{a=0}]$。

\section{可经验验证的中介}

$\E\big[Y^{a=1,M^{a=0}}\big]$ 的干预主义解释，只有在关于 $N$ 与 $O$ 可分离效应的故事正确、且图 23.3 表示 FFRCISTG 时才成立。干预主义解释的一个优点是：这个故事的正确性原则上可以通过随机试验被经验地反驳，下面就来描述。

一年后无尼古丁香烟上市时，我们做一项新的随机试验：把吸烟者的随机样本随机分配到三组之一：戒烟（$A=N=O=0$）、继续吸标准香烟（$A=N=O=1$）、或继续吸无尼古丁香烟（$N=0,O=1$）。也就是说，新试验有与原试验相同的两个组，另加一个无尼古丁香烟组。若无时间趋势，前两组的平均结局在这两项试验中应当相同。我们假定两项试验的样本量都足够大，可以忽略抽样变异。

由随机化，新试验中的平均结局 $\E[Y\given N=0,O=1]$ 期望等于 $\E\big[Y^{n=0,o=1}\big]$。因此，若我们的故事正确，我们预期 $\E[Y\given N=0,O=1]$ 将等于早先试验的中介公式。否则，$N$ 与 $O$ 的故事就被经验反驳了，这意味着下述假定中至少有一个不正确：(i) 尼古丁对结局没有直接效应；(ii) 其他非尼古丁成分对中介没有直接效应；(iii) 图 23.2 的因果图上 $M$ 与 $Y$ 没有未测量的共同原因 $U$。

若 $\E[Y\given N=0,O=1]$ 与中介公式不同，我们可以用新试验的数据考察假定 (i)—(iii) 中哪些不成立。为此，先检查在 $N=0$ 者中 $O$ 与 $M$ 是否有关联，即检查下式是否成立：
\[
\E[M\given N=0,O=1]-\E[M\given N=0,O=0]\neq 0.
\]

若如此，则假定 (ii) 被反驳，应把箭头 $O\rightarrow M$ 加入图 23.2（并把相应的箭头加入图 22.3）。

接着检查在 $M$ 与 $O$ 的联合水平内 $N$ 与 $Y$ 是否有关联，即对某些 $m$ 值是否有
\[
\E[Y\given N=1,O=1,M=m]-\E[Y\given N=0,O=1,M=m]\neq 0.
\]
若如此，则假定 (i) 被反驳（这意味着应把箭头 $N\rightarrow Y$ 加入因果图），或者存在 $M$ 与 $Y$ 的未测量共同原因 $U$（应把它加入因果图）。见精要 23.1。

\begin{finepoint}[可分离效应诸假定的经验否证]
\setlength{\parindent}{2em}
假定我们做了对 $N$ 与 $O$ 施加干预的三组试验，结果发现在 $M$ 与 $O$ 的联合水平内 $N$ 与 $Y$ 有关联。那么，如正文所述，我们断定假定 (i) 或假定 (iii) 之一不成立，或二者皆不成立。

要确定这两个假定中哪一个不成立，我们需要另一项随机试验，比如说有 8 个组，在其中同时对 $M$ 以及 $N$、$O$ 施加干预。若 $N$ 除经由 $M$ 之外对 $Y$ 没有直接效应，则在八组试验中给定 $M$ 与 $O$ 时 $N$ 将与 $Y$ 独立。若 $M$ 与 $Y$ 没有未测量的共同原因，则在 $N$ 与 $O$ 的各个水平内，给定 $M$ 时 $Y$ 的条件分布在三组试验与八组试验中应当相同。

有人会问：如果我们正确地假定了图 23.1 表示一个 $M$ 与 $Y$ 没有共同原因的 FFRCISTG 模型，图 23.2 怎么还会遗漏 $M$ 与 $Y$ 的未测量共同原因？解释是：共同原因可能只在给每个人指派不同 $N$ 与 $O$ 取值的干预下才存在（即才被激活）。与图 23.1 相联系的那个 FFRCISTG 模型，只考虑给每个人指派相同取值的干预 $A=N=O=1$ 与 $A=N=O=0$，而图 23.2 所联系的模型则不然。有趣的是，若 $Y$ 与 $M$ 的共同原因 $U$ 是新试验中 $\E[Y\given N=0,O=1]$ 不同于中介公式的唯一原因，则对每个个体仍有 $Y^{n=0,o=1}=Y^{a=1,M^{a=0}}$。然而，$\E\big[Y^{a=1,M^{a=0}}\big]=\E\big[Y^{n=0,o=1}\big]$ 不能由中介公式识别。Robins 等 (2022) 描述了一项关于河盲症治疗的假想但现实的研究，其中图 23.1 是 FFRCISTG，而图 23.2 不是，因为后者遗漏了 $Y$ 与 $M$ 的共同原因 $U$。
\end{finepoint}

现在设想两年后新试验结果出来：对应于两项试验中同一干预的组，平均结局估计相等，但新试验第三组的平均结局估计 $\E[Y\given N=0,O=1]$ 与第一项试验的中介公式不同。不同的人会如何反应？

无论是曾假定图 23.1 为 NPSEM-IE 的人，还是只假定它为 FFRCISTG 的人，都会同意：$O$ 与 $N$ 的可分离效应这个故事不正确，$\E\big[Y^{n=0,o=1}\big]$ 应由新试验第三组的平均结局正确估计，而不是由早先试验的中介公式估计。

曾假定 NPSEM-IE 的人可以继续相信均值 $\E\big[Y^{a=1,M^{a=0}}\big]$ 仍等于中介公式，但他们无法再把纯直接效应作为“改吸无尼古丁香烟的效应”来论证其政策意义。曾假定 FFRCISTG 的人可能对原来的中介公式兴趣不大。相反，他们会对自己从三组随机试验中学到的、关于无尼古丁香烟对结局之效应的知识感兴趣。他们也会乐于知道下述假定中有一个或多个不成立：(i) 尼古丁对结局没有直接效应；(ii) 其他非尼古丁成分对中介没有直接效应；(iii) $M$ 与 $Y$ 没有未测量的共同原因 $U$。

继续假定图 23.1 表示 FFRCISTG 模型，且未来无尼古丁香烟试验中的 $\E\big[Y^{n=0,o=1}\big]$ 与中介公式不同。一个有趣的问题是：\marginnote{\textcolor{black!85}{若 $\E\big[Y^{a=1,M^{a=0}}\big]$ 总能由中介公式点识别，Robins 与 Richardson (2010) 的锐界就会被违反。}}处理 $A$ 是否总能分解为其他（也许是未知的）可干预成分 $N'$ 与 $O'$，使得图 23.2 是 FFRCISTG 且无直接效应假定 (i) 与 (ii) 成立？若如此，$\E\big[Y^{a=1,M^{a=0}}\big]$ 就会等于 $\E\big[Y^{n'=0,o'=1}\big]$，二者都等于中介公式。答案是否定的，因为否则 $\E\big[Y^{a=1,M^{a=0}}\big]$ 就会总能由中介公式点识别，而事实并非如此。

\section{干预主义的中介理论}

本章介绍了一种干预主义的因果中介理论，它有别于基于纯直接效应与总间接效应的标准路径。由于干预主义理论是为了回应标准理论中公认的缺陷而发展的，我们先描述了标准理论及其问题所在——这些问题立足于跨世界（嵌套）反事实的使用。然而，干预主义理论可以视为自足的，它提供了一个自洽的框架来讨论中介，而不必诉诸跨世界嵌套反事实。\mnote{这种干预主义的中介理论由 Robins 与 Richardson (2010) 提出，并由 Robins、Richardson 与 Shpitzer (2022) 扩展。该理论已被推广到生存分析（Dídelez Aalen 等 2020）、竞争风险（Stensrud 等 2020, 2021）以及存在干扰的情形（Shpitzer 等 2021）。}

为此，我们用一个（简化的）随机试验作为干预，其中处理是时点固定的。关键想法是把中介问题重新表述为一个关于“对处理的实质上可分离成分（$N$ 与 $O$）的干预”之效应的问题。若 $N$ 与 $O$ 是 $A$ 的可分离效应，那么在一项有六个组的新随机试验中：

\begin{enumerate}
\item 用 $a=1$ 处理
\item 用 $a=0$ 处理
\item 用 $n=1,o=1$ 处理
\item 用 $n=0,o=0$ 处理
\item 用 $n=0,o=1$ 处理
\item 用 $n=1,o=0$ 处理
\end{enumerate}

则下述两个陈述为真。第一，第一组与第三组中结局 $Y$ 的均值相等，第二组与第四组中结局 $Y$ 的均值也相等。第二，$\E\big[Y^{n=0,o=1}\big]$ 被识别为第五组的平均结局，$\E\big[Y^{n=1,o=0}\big]$ 被识别为第六组的平均结局——无论 $\E\big[Y^{n=0,o=1}\big]$ 是否能由观测数据识别。\mnote{Stensrud 等 (2021) 强调了上述“六组试验各组平均结局相等”的必要性。}

\begin{figure}[htbp]
\centering
\begin{tikzpicture}[causal]
  \node[vtx] (A) at (0,0) {$A$};
  \node[vtx] (H) at (1.8,0) {$H$};
  \node[vtx] (M) at (1.8,1.2) {$M$};
  \node[vtx] (Y) at (3.8,0) {$Y$};
  \draw[ed] (A) -- (H);
  \draw[ed] (H) -- (M);
  \draw[ed] (H) -- (Y);
  \draw[ed] (A) to[bend right=40] (Y);
\end{tikzpicture}
\caption{}
\label{fig:23.4}
\end{figure}

\begin{figure}[htbp]
\centering
\begin{tikzpicture}[causal]
  \node[vtx] (A) at (0,0) {$A$};
  \node[vtx] (N) at (1.7,0.9) {$N$};
  \node[vtx] (O) at (1.7,-0.9) {$O$};
  \node[vtx] (H) at (3.5,0) {$H$};
  \node[vtx] (M) at (3.5,1.2) {$M$};
  \node[vtx] (Y) at (5.3,0) {$Y$};
  \draw[ed] (A) -- (N);
  \draw[ed] (A) -- (O);
  \draw[ed] (N) -- (H);
  \draw[ed] (O) -- (Y);
  \draw[ed] (H) -- (M);
  \draw[ed] (H) -- (Y);
\end{tikzpicture}
\caption{}
\label{fig:23.5}
\end{figure}

我们可以用当前可得的 $A$、$M$、$Y$ 数据，在如下假定下识别 $N$ 与 $O$ 的可分离效应：(i) 中介 $M$ 与结局 $Y$ 没有未测量的共同原因；(ii) 成分 $O$ 对中介 $M$、成分 $N$ 对结局 $Y$ 都没有直接效应。特别地，$\E\big[Y^{n=0,o=1}\big]$ 等于中介公式。

\begin{finepoint}[中介为替代变量时的可分离效应]
\setlength{\parindent}{2em}
下面我们采用第 9.5 节引入的因果图理论。由于 $M$ 到 $Y$ 的箭头只有在“$M$ 对 $Y$ 的效应存在定义良好的干预”时才有因果解释，所以当这些干预不存在时，图 23.1 的 DAG 不是因果的。也就是说，$M$ 不是中介，而是某个未知变量 $H$ 的替代；$H$ 才是存在定义良好干预的真正中介。图 23.4 的因果 DAG 刻画了这一情形。

考虑图 23.5 的因果图，它是在图 23.4 上加入可分离成分 $N$ 与 $O$ 的扩展。与图 23.2 不同，在图 23.5 中给定 $M$ 与 $O$ 时 $N$ 与 $Y$ 不是 d-分离的。设想我们惊讶地发现，在三组试验数据中给定 $M$ 与 $O$ 时 $N$ 与 $Y$ 独立。应当如何解释这一观察？有 4 种可能：a) 我们弄错了，图 23.1 而非 23.4 才是真实因果图，$M$ 是真正的因果中介；b) $M$ 是真正因果中介 $H$ 的一一确定性函数；c) 图 23.4 中存在非确定性的忠实性违背；d) $N$ 与 $O$ 并非条件独立，只是三组试验样本太小，未能检出其相依性。因果发现的支持者若见三组试验样本量很大，往往倾向于采用解释 a)（见技术注记 10.7）。
\end{finepoint}

下面再讨论一个宁愿采用干预主义路径、而不采用早期因果中介路径的理由。具体而言，对假定的中介 $M$ 的干预常常定义不良，此时像 $Y^{a,m}$ 这样的反事实就没有意义。于是，无论是纯直接效应还是基于 $M$ 的受控直接效应（见技术注记 22.1）都不存在。相反，干预主义理论关心的是确实存在的效应：即使对 $M$ 的干预定义不良，只要存在实质上有意义、可按上述方式干预的可分离成分 $N$ 与 $O$ 即可。

总之，干预主义的中介路径包含以下成分。

\begin{itemize}
\item 假定处理 $A$ 可分解为多个实质上有意义的可分离成分，每个成分都对处理的总体效应有贡献，且原则上可被独立干预；
\item 识别可分离成分之效应所需的假定，原则上可在未来真正干预这些成分的随机试验中得到经验验证；
\item 对所谓中介的良好定义干预不必存在；
\item 由于可分离成分具有实质特定性，可与领域专家进行更有效的沟通；
\item 当识别假定成立时，用于识别的 g 公式与中介公式相同；但所识别的因果效应指的是对可分离成分的干预。
\end{itemize}

为简单起见，我们只考虑了两个可分离成分的例子，但干预主义框架可以容纳处理的多个可分离成分，包括随时间变化的成分。

\begin{figure}[htbp]
\centering
\begin{tikzpicture}[causal]
  \node[vtx] (H) at (0,-1.3) {$H$};
  \node[vtx] (L) at (0,0) {$L$};
  \node[vtx] (A) at (1.9,0) {$A$};
  \node[vtx] (Y) at (3.8,0) {$Y$};
  \draw[ed] (H) -- (L);
  \draw[ed] (L) -- (A);
  \draw[ed] (A) -- (Y);
  \draw[ed] (L) to[bend left=42] (Y);
  \draw[ed] (H) -- (Y);
\end{tikzpicture}
\caption{}
\label{fig:23.6}
\end{figure}

\begin{figure}[htbp]
\centering
\begin{tikzpicture}[causal]
  \node[vtx] (H) at (0,-1.4) {$H$};
  \node[vtx] (L) at (0,0) {$L$};
  \node[vtx] (N) at (1.7,0.9) {$N$};
  \node[vtx] (O) at (1.7,-0.9) {$O$};
  \node[vtx] (A) at (3.5,0.9) {$A$};
  \node[vtx] (Y) at (5.2,0) {$Y$};
  \draw[ed] (H) -- (L);
  \draw[edb] (L) -- (N);
  \draw[edb] (L) -- (O);
  \draw[ed] (N) -- (A);
  \draw[ed] (A) -- (Y);
  \draw[ed] (O) -- (Y);
\end{tikzpicture}
\caption{}
\label{fig:23.7}
\end{figure}

\begin{figure}[htbp]
\centering
\begin{tikzpicture}[causal,swig/.style={draw,rounded corners=2pt,inner sep=2.5pt,minimum height=6mm,font=\small}]
  \node[vtx] (H) at (0,-1.4) {$H$};
  \node[vtx] (L) at (0,0) {$L$};
  \node[swig] (N) at (2.0,0.9) {$N\,|\,n$};
  \node[vtx] (O) at (2.0,-0.9) {$O$};
  \node[swig] (A) at (4.1,0.9) {$A^{n}$};
  \node[swig] (Y) at (5.9,0) {$Y^{n}$};
  \draw[ed] (H) -- (L);
  \draw[edb] (L) -- (N);
  \draw[edb] (L) -- (O);
  \draw[ed] (N) -- (A);
  \draw[ed] (A) -- (Y);
  \draw[ed] (O) -- (Y);
\end{tikzpicture}
\caption{}
\label{fig:23.8}
\end{figure}

本章回顾了中介的理论概念，但没有强调估计纯直接效应与可分离直接效应的数据分析的实际方面。上一章在介绍受控直接效应与主分层直接效应时采用了类似的做法。我们希望这里的讲述已阐明：干预主义路径在识别直接效应方面具有科学优势。

\begin{techpoint}[路径特定效应与前门公式]
\setlength{\parindent}{2em}
在图 23.6（图 9.14 的修改版）下，$L$ 对 $Y$ 的总效应不可识别。这是因为通路 $L\rightarrow Y$ 对总效应的贡献无法与 $H$ 的混杂效应分开。相反，在缺少 $L\rightarrow Y$ 边的图 9.14 原图中，我们在精要 9.5 中已证明 $L$ 对 $Y$ 的总效应可由前门公式识别，因为 $L\rightarrow A\rightarrow Y$ 是唯一的因果通路。

现在考虑：当 $L\rightarrow Y$ 边存在时，$L$ 沿通路 $L\rightarrow A\rightarrow Y$ 的效应是否也能被识别，且仍由同一个前门公式识别。本着本章的精神，我们用精要 9.5 中的实质故事为这一问题提供干预主义的表述。图 23.7 是一个扩展图，其中 $N$ 记录向负责开具药物 $A$ 的医生报告的 BMI 值 $L$，$O$ 记录用于决定是否转诊接受物理治疗与饮食咨询的 $L$ 值。粗箭头表示确定性关系，因为在实际观测数据中我们总有 $L=N=O$。图 23.8 是一个 SWIG，表示一项（不合伦理的）干预：向医生报告一个不同于真实值 $L=N$ 的 BMI 值 $n$，但转诊仍使用真实值 $L=O$。与本章前面的做法不同，这里只对 $N$ 施加了干预。此外，正如我们的目标所要求，对 $n$ 的干预对 $Y$ 的效应被限制在经由 $A$ 的通路上。由于即使在干预后的世界里仍有 $L\equiv O$，我们不需要 $O$。于是我们从图 23.7 与 23.8 中删去 $O$，$Y$ 又成为 $L$ 的子节点。此时只需注意到在图 23.8 上平凡地有 $Y^{n}\indep N\given L$，$\E[Y^{n}]$ 就应由基于图 23.7 的 g 公式给出：
\begin{align*}
&\sum_{l,a} \E[Y\given A=a,L=l]\,\Pr[L=l]\,\Pr[A=a\given N=n]\\
&\qquad= \sum_a \left\{\sum_l \E[Y\given A=a,L=l]\,\Pr[L=l]\right\}\Pr[A=a\given N=n],
\end{align*}
其中最后一个等号依据数据中的确定性 $L\equiv N$。右端确实就是前门公式。

然而，这一推导由于确定性带来的零测集而略显启发式。把确定性与技术注记 21.12 中证明前门公式所用的方法结合起来，即可容易地得到严格证明。还要注意，从实质上看，人们一般不会把 $N$ 与 $O$ 视为 $L$ 的可分离成分。但这无关紧要。要紧的是：我们的实质因果故事蕴含一个扩展图 23.7，其中含有与实际世界中的 $L$ 确定性相关的新干预变量 $N$（Stensrud 等，2023）。Wen 等 (2023) 独立且同时得到了与这一技术注记本质等价的结果。Fulcher 等 (2020) 此前已证明，在与图 23.6 相联系的 NPSEM-IE 模型下，上述前门公式识别跨世界反事实量 $\E\big[Y^{L,A^{l=n}}\big]$。Wen 等的工作与本技术注记都可看作对 Fulcher 等结果所作的干预主义重述。
\end{techpoint}

在缺乏对中介（就受控直接效应而言）或处理成分（就可分离直接效应而言）真正施加干预的随机试验时，所有这些方法都依赖观测数据，因而依赖本书通篇讨论的可交换性假定。因此，有效的中介分析除了本章讨论的其他假定外，还需要对处理效应的混杂因子与中介效应的混杂因子都做调整。图 23.1 的因果图只是用来在简化情境中探讨中介相关理论问题的教学工具，并不是对多数因果中介研究的现实刻画。在实践中，因果中介分析是观测分析，依赖的假定比非中介分析更为大胆。\mnote{对可分离效应，若存在 $M$ 与 $Y$ 的已测量共同原因、而它又是 $N$ 与 $O$ 的子代，则还会出现额外问题（Robins 与 Richardson 2010；Robins 等 2022）。}
